\subsection{Second Order Systems}
\label{sec:theory_cs_second_order}

A second order system is one that can be described by a
differential equation of order 2 in time domain which
transforms into polynomial with degree 2 in Frequency
domain. The general form of a second order transfer
function is given by:
\begin{align*}
  H(s) = \frac{V_{out}}{V_{in}} = \frac{\text{Zeros}}{\text{Poles}} = \frac{N(s)}{D(s)} = \frac{b_2 s^2 + b_1 s + b_0}{a_2 s^2 + a_1 s + a_0} = \underbrace{\frac{b_2 s^2}{a_2 s^2+a_1 s+a_0}}_{\text{Highpass}} + \underbrace{\frac{b_1 s}{a_2 s^2+a_1 s+a_0}}_{\text{Bandpass}} + \underbrace{\frac{b_0}{a_2 s^2+a_1 s+a_0}}_{\text{Lowpass}}
\end{align*}

\paragraph{Mathematical Overview of Poles:}
There are multiple ways to represent a second order
polynomial. I have listed four common forms below:
\begin{enumerate}
  \item The Mathematical form: This is the general representation
    of a second order polynomial.
    \begin{equation*}
      a x^2 + b x + c
    \end{equation*}
  \item The Normalized form: We divide the entire polynomial
    by the coefficient of \( s^2 \) (= \( x^2 \)) to make
    the leading coefficient 1.
    \begin{equation*}
      s^2 + \frac{a_1}{a_2}s + \frac{a_0}{a_2}
    \end{equation*}
  \item The Factored form: We express the polynomial
    in terms of its roots (poles).
    \begin{align*}
      (s + p_1)(s + p_2) \\
      s^2 + (p_1 + p_2)s + p_1 p_2
    \end{align*}
    where \( p_1 \) and \( p_2 \) are the roots (poles) of the polynomial.
  \item The Natural form: We express the polynomial
    in terms of its natural frequency and damping ratio.
    \begin{align*}
      s^2 + 2 \zeta \omega_n s + \omega_n^2 \\
      s^2 + \frac{\omega_n}{Q} s + \omega_n^2
    \end{align*}
    where \( \omega_n \) is the natural frequency and \( \zeta \) is the
    damping ratio and \( Q\) is the quality factor.
\end{enumerate}

Here is a summary table of the coefficients across different forms:

\begin{center}
  \begin{tabularx}{\linewidth}{|X|X|X|X|X|}
    \hline
    \textbf{Coefficient} & \textbf{Mathematical} & \textbf{Normalized} & \textbf{Factored} & \textbf{Natural} \\
    \hline
    \(s^2\) & \(a\) & 1 & 1 & 1 \\
    \hline
    \(s\) & \(b\) & \(\displaystyle\frac{a_1}{a_2}\) & \(p_1 + p_2\) & \(2\zeta\omega_n\) \\
    \hline
    Constant & \(c\) & \(\displaystyle\frac{a_0}{a_2}\) & \(p_1 p_2\) & \(\omega_n^2\) \\
    \hline
  \end{tabularx}
\end{center}

Next, we explore the relationships between coefficients, poles,
natural frequency, and damping ratio through three key comparisons.
\begin{enumerate}
  \item \textbf{Normalized Form vs. Factored Form}
    \noindent From the table, we have:
    \[ \frac{a_1}{a_2} = p_1 + p_2 \quad \text{and} \quad \frac{a_0}{a_2} = p_1 p_2 \]
    The discriminant
    \(\displaystyle \Delta = \left( \frac{a_1}{a_2} \right)^2 - 4 \left( \frac{a_0}{a_2} \right) \)
    determines three pole configurations:
    \begin{enumerate}
      \item \textbf{Two distinct real poles (overdamped):} Assume \( p_2 = k p_1 \) with \( k \in \mathbb{R} \), \( k \neq 1 \).
        \begin{align*}
          \frac{a_1}{a_2} &= p_1(1 + k) \quad \Rightarrow \quad p_1 = \frac{a_1}{a_2(1 + k)} \tag{i} \\
          \frac{a_0}{a_2} &= k p_1^2 \quad \Rightarrow \quad k p_1^2 = \frac{a_0}{a_2} \tag{ii}
        \end{align*}
        Substituting (i) into (ii):
        \begin{align*}
        k \left( \frac{a_1}{a_2(1 + k)} \right)^2 = \frac{a_0}{a_2}
        \quad \Rightarrow \quad
        \frac{k}{(1+k)^2} = \frac{a_0 a_2}{a_1^2}
        \end{align*}
        For large \(k\):
        \[ k \approx \frac{a_1^2}{a_0 a_2} \]
      \item \textbf{One repeated real pole (critically damped):} Let \( p_1 = p_2 = p \).
        \begin{align*}
          \frac{a_1}{a_2} &= 2p \quad \Rightarrow \quad p = \frac{a_1}{2a_2} \tag{i} \\
          \frac{a_0}{a_2} &= p^2 \quad \Rightarrow \quad p^2 = \frac{a_0}{a_2} \tag{ii}
        \end{align*}
        Substituting (i) into (ii):
        \begin{align*}
          \left( \frac{a_1}{2a_2} \right)^2 = \frac{a_0}{a_2}
          \quad \Rightarrow \quad
          a_1^2 = 4 a_0 a_2
        \end{align*}
      \item \textbf{Two complex conjugate poles (underdamped):} Let \( p_{1,2} = \alpha \pm j\beta \), with \( \alpha, \beta \in \mathbb{R} \).
        \begin{align*}
          \frac{a_1}{a_2} &= 2\alpha \quad \Rightarrow \quad \alpha = \frac{a_1}{2a_2} \tag{i} \\
          \frac{a_0}{a_2} &= \alpha^2 + \beta^2 \quad \Rightarrow \quad \alpha^2 + \beta^2 = \frac{a_0}{a_2} \tag{ii}
        \end{align*}
        Substituting (i) into (ii):
        \begin{align*}
          \left( \frac{a_1}{2a_2} \right)^2 + \beta^2 = \frac{a_0}{a_2}
          \quad \Rightarrow \quad
          \beta = \sqrt{\frac{a_0}{a_2} - \frac{a_1^2}{4a_2^2}}
        \end{align*}
    \end{enumerate}
    These relations help determine pole locations from
    polynomial coefficients and help in pole placement
    design (chose k to set pole separation).
  \item \textbf{Normalized Form vs. Natural Form}
    \noindent From the table:
    \[ \frac{a_1}{a_2} = 2\zeta\omega_n \quad \text{and} \quad \frac{a_0}{a_2} = \omega_n^2 \]
    From the second equation:
    \[ \omega_n = \sqrt{\frac{a_0}{a_2}} \tag{ii} \]
    Substituting (ii) into the first:
    \begin{align*}
    \frac{a_1}{a_2} = 2\zeta \sqrt{\frac{a_0}{a_2}}
    \quad \Rightarrow \quad
    \zeta = \frac{a_1}{2\sqrt{a_0 a_2}}
    \end{align*}
    Since quality factor \( Q = 1/(2\zeta) \):
    \[ Q = \frac{\sqrt{a_0 a_2}}{a_1} \]
    Unlike the factored form, comparison with the natural form does
    not require separate algebraic cases; the three damping regimes
    emerge solely from the numerical value of \(\zeta\). These
    relations help determine natural frequency and damping ratio from
    polynomial coefficients.
  \item \textbf{Factored Form vs. Natural Form}
    \noindent From the table:
    \[ p_1 + p_2 = 2\zeta\omega_n \quad \text{and} \quad p_1 p_2 = \omega_n^2 \]
    From the second equation:
    \[ \omega_n = \sqrt{p_1 p_2} \tag{ii} \]
    Substituting (ii) into the first:
    \begin{align*}
      p_1 + p_2 = 2\zeta \sqrt{p_1 p_2}
      \quad \Rightarrow \quad
      \zeta = \frac{p_1 + p_2}{2\sqrt{p_1 p_2}}
    \end{align*}
    Expressed in terms of quality factor:
    \[ Q = \frac{\sqrt{p_1 p_2}}{p_1 + p_2}\]
    We now examine the three pole configurations:
    \begin{enumerate}
      \item \textbf{Two distinct real poles (overdamped):} With \( p_2 = k p_1 \),
        \[ Q = \frac{\sqrt{k p_1^2}}{p_1 + k p_1} = \frac{\sqrt{k}}{1+k} \]
        For large k:
        \[ Q \approx \frac{1}{\sqrt{k}} \]
      \item \textbf{One repeated real pole (critically damped):} With \( p_1 = p_2 = p \),
        \[ Q = \frac{2\sqrt{p^2}}{p + p} = 1 \]
      \item \textbf{Two complex conjugate poles (underdamped):} With \( p_{1,2} = \alpha \pm j\beta \),
        \begin{align*}
           Q = \frac{2\sqrt{(\alpha + j\beta)(\alpha - j\beta)}}{2\alpha}
           = \frac{2\sqrt{\alpha^2 + \beta^2}}{2\alpha} \\
            Q = \frac{\sqrt{\alpha^2 + \beta^2}}{\alpha}
         \end{align*}
    \end{enumerate}
    These relations help determine natural frequency and
    damping ratio from pole locations.
\end{enumerate}
Here's  table summarizing the relationships:
\begin{center}
  \begin{tabularx}{\linewidth}{|X|X|X|X|} \hline
    \textbf{Comparison} & \textbf{Key Equations} & \textbf{Pole Configurations} & \textbf{Design Insights} \\ \hline
    Normalized vs. Factored & \(\frac{a_1}{a_2} = p_1 + p_2\) \newline \(\frac{a_0}{a_2} = p_1 p_2\) & Over: \(k \approx \frac{a_1^2}{a_0 a_2}\) \newline Critical: \(a_1^2 = 4 a_0 a_2\) \newline Under: \(\beta = \sqrt{\frac{a_0}{a_2} - \frac{a_1^2}{4a_2^2}}\) & Determine pole locations from coefficients; pole placement design. \\ \hline
    Normalized vs. Natural & \(\frac{a_1}{a_2} = 2\zeta\omega_n\) \newline \(\frac{a_0}{a_2} = \omega_n^2\) & \(\omega_n = \sqrt{\frac{a_0}{a_2}}\) \newline \(\zeta = \frac{a_1}{2\sqrt{a_0 a_2}}\) & Determine natural frequency and damping ratio from coefficients. \\ \hline
    Factored vs. Natural & \(p_1 + p_2 = 2\zeta\omega_n\) \newline \(p_1 p_2 = \omega_n^2\) & Over: \(Q \approx \frac{1}{\sqrt{k}}\) \newline Critical: \(Q = 1\) \newline Under: \(Q = \frac{\sqrt{\alpha^2 + \beta^2}}{\alpha}\) & Determine natural frequency and damping ratio from pole locations. \\ \hline
  \end{tabularx}
\end{center}

\paragraph{Example:}
\begin{enumerate}
  \item Given the second order polynomial \( 3s^2 + 12s + 27 \),
  determine the pole configuration, natural frequency, and
  damping ratio.
  \paragraph{Solution:}
  \begin{enumerate}
    \item Identify coefficients:
      \[ a_2 = 3, \quad a_1 = 12, \quad a_0 = 27 \]
    \item Calculate natural frequency:
      \begin{align*}
        \omega_n = \sqrt{\frac{a_0}{a_2}} = \sqrt{\frac{27}{3}} = \sqrt{9} = 3 \text{ rad/s}
      \end{align*}
    \item Calculate damping ratio:
      \begin{align*}
        \zeta = \frac{a_1}{2\sqrt{a_0 a_2}} = \frac{12}{2\sqrt{27 \cdot 3}} = \frac{12}{2\sqrt{81}} = \frac{12}{18} = \frac{2}{3} \approx 0.6667
      \end{align*}
    \item Determine pole configuration:
      Since \( \zeta \approx 0.6667 < 1 \), the system is underdamped with
      complex conjugate poles.
  \end{enumerate}
  \item Given the second order polynomial \( 2s^2 + (8 + x)s + 16 \),
  determine the value of \( x \) for critical damping.
  \paragraph{Solution:}
  \begin{enumerate}
    \item Identify coefficients:
      \[ a_2 = 2, \quad a_1 = 8 + x, \quad a_0 = 16 \]
    \item Apply critical damping condition:
      \begin{align*}
        a_1^2 = 4 a_0 a_2 \\
        (8 + x)^2 = 4 \cdot 16 \cdot 2 \\
        (8 + x)^2 = 128 \\
        8 + x = \sqrt{128} \\
        x = \sqrt{128} - 8 \approx 3.3137
      \end{align*}
  \end{enumerate}
  Thus, \( x \approx 3.3137 \) for critical damping.
\end{enumerate}
